\section{Conclusion}
This paper presented an unsupervised dual-stream masked autoencoder for intrusion detection in Automotive Ethernet networks, addressing two fundamental challenges in automotive IDS development: protocol heterogeneity and the scarcity of annotated attack traffic. A sliding window of 64 consecutive packets is encoded through two complementary streams -- a Transformer encoder that captures global inter-packet dependencies from raw byte content, conditioned on a time-aware positional encoding derived from real inter-packet capture timing rather than fixed sequence rank, and a residual 1D-CNN encoder that models local per-byte temporal dynamics from inter-packet byte differences -- unified through a delta-conditioned gated residual fusion mechanism. The model is trained via packet-level masked reconstruction on unlabeled normal traffic alone, and anomalies are flagged at inference time by thresholding the mean reconstruction error against a percentile-based threshold derived from normal validation traffic.

Experiments on the CarDS Automotive Ethernet dataset, spanning 27 real-vehicle attack files across 23 attack configurations, show that the proposed method achieves an overall F1-score of 0.9808, AUROC of 0.9974, and AUPRC of 0.9972 at its selected operating threshold ($P_{99.9}$), essentially matching a fully supervised baseline (AUROC 0.999) while requiring no attack labels. Detection is near-perfect on Nmap-based scanning and spoofing attacks (F1 $>$ 0.95) but comparatively harder on RTP stream replacement (avg.\ F1 = 0.8709), which alters only payload content while preserving packet-level structure and timing. Ablation studies confirm that the dual-stream architecture and the time-aware positional encoding each contribute independently to detection performance, with the latter also substantially widening the range of thresholds over which performance remains stable.

Several limitations point to directions for future work. The method was evaluated on a single real-vehicle dataset covering two attack families (RTP stream replacement and Nmap-based scanning/spoofing); broader validation across vehicle platforms, network topologies, and attack types is needed to establish generalizability. Threshold calibration, the one place labeled attack data enters the pipeline, currently relies on a percentile search over labeled validation attacks, and less threshold-sensitive or fully label-free scoring schemes merit further study. The mechanism by which time-aware positional encoding improves RTP detection despite the attack itself preserving inter-packet timing also remains unclear and warrants attention-weight or controlled perturbation analysis. Finally, deploying the framework in online, resource-constrained settings -- including latency characterization and model compression for embedded ECU targets -- is an important direction for practical automotive IDS.